% Clase
\documentclass[11pt,a4paper,spanish,twoside]{book}

% Órdenes auxiliares
% Español
\usepackage[spanish]{babel}
\usepackage[utf8]{inputenc}
\usepackage[T1]{fontenc}
\usepackage{lmodern}

% Imágenes
\usepackage[pdftex]{graphicx}
\usepackage{latexsym}
\usepackage{fancybox}

% Ruta para las imágenes
\graphicspath{{img/}}

% Rotaciones
\usepackage[twoside]{rotating}

% Multirow
\usepackage{multirow}

% Referencias
\usepackage[spanish]{varioref}
\usepackage[pdftex,colorlinks=true,linkcolor=black]{hyperref}

% Colores
\usepackage{color}
\usepackage{colortbl}

% Párrafos
\setlength{\parskip}{6pt}

% Code for creating empty pages
% No headers on empty pages before new chapter
\makeatletter
\def\cleardoublepage{\clearpage\if@twoside \ifodd\c@page\else
    \hbox{}
    \thispagestyle{empty}
    \newpage
    \if@twocolumn\hbox{}\newpage\fi\fi\fi}
\makeatother \clearpage{\pagestyle{empty}\cleardoublepage}

\usepackage{titlesec}

\titleformat{\chapter}[display]
  {\bfseries\Large}
  {\filleft\MakeUppercase{\chaptertitlename} \Huge\thechapter}
  {4ex}
  {\titlerule
   \vspace{2ex}%
   \filright}
  [\vspace{2ex}%
   \titlerule]

% Fancyhdr - Encabezados y pies de página
\usepackage{fancyhdr}
% Márgenes
\headsep=8mm
\footskip=14mm

% Fancy Header Style Options
\pagestyle{fancy}               % Sets fancy header and footer
\fancyfoot{}                    % Delete current footer settings

% Sin mayúsculas en la cabecera
\lhead{\nouppercase{\rightmark}}
\rhead{\nouppercase{\leftmark}}

% Capítulo
\renewcommand{\chaptermark}[1]{ % Lower Case Chapter marker style
  \markboth{\chaptername\ \thechapter.\ #1}{}} 

% Sección
\renewcommand{\sectionmark}[1]{ % Lower case Section marker style
  \markright{\thesection.\ #1}} 

% Página
\fancyhead[LE,RO]{\bfseries\thepage} % Page number (boldface) in left on even
                                     % pages and right on odd pages

% ------------------ Macro para encabezados y pies de página-------------------
%    Uso: \encabezado{pares(pag izquierda)}
% -----------------------------------------------------------------------------
\def\encabezado{
  \fancyhead[RE]{\bfseries\leftmark}     % In the right on even pages
  \fancyhead[LO]{\bfseries\rightmark}  % In the left on odd pages
  \renewcommand{\headrulewidth}{0.5pt} % Width of head rule
}
% -----------------------------------------------------------------------------

% ------------------ Macro para insertar una imagen ---------------------------
%    Uso: \imagen{nombreFichero}{Ancho(cm))}{Etiqueta}{Identificador}
% -----------------------------------------------------------------------------
\usepackage{float}
\usepackage{ifthen}
\def\imagen#1#2#3#4{
  \begin{figure}[H]
    \begin{center}
      \ifthenelse{\equal{#2}{}}
      {\includegraphics{#1}}{\resizebox{#2cm}{!}{\includegraphics{#1}}}
      \ifthenelse{\equal{#3}{}}{}{\caption{#3}}
      \label{#4}
    \end{center}
  \end{figure}
}
% -----------------------------------------------------------------------------


% ------------------ Macro para la portada ------------------------------------
%    Uso: \portada{asignatura}{titulo}{subtítulo}{autor}{fecha}
% -----------------------------------------------------------------------------
\def\portada#1#2#3#4#5{
  \thispagestyle{empty}
  \vspace*{-3.3cm}

  \begin{minipage}[t]{14cm}
    \begin{center}
      \includegraphics[scale=0.25]{logoesi}
  
      \vspace*{1.5cm}
      {\Large \textbf{Universidad de Castilla-La Mancha\\ 
          Escuela Superior de Informática}\\}
    
      \vspace*{1.2cm}
      {\huge \textbf{#1}\\}
    
      \vspace*{1.5cm}
      {\huge #2}\\{\Large #3\\}
    
      \vspace*{1.5cm}
      {\large #4\\}
      \vspace*{1.4cm}
      \large{#5}
    \end{center}
  \end{minipage}

  \newpage
%   \vspace*{1cm}
%   \thispagestyle{empty} 
%   \newpage
}
% -----------------------------------------------------------------------------


% ------------------ Macro para la licencia -----------------------------------
%    Uso: \portada{asignatura}{titulo}{subtítulo}{autor}{fecha}
% -----------------------------------------------------------------------------
\def\licencia#1{
  \thispagestyle{empty}  % Suprime la numeración de esta página
  \vspace*{16.5cm}
  \begin{small}
    \copyright~ #1. Se permite la copia, distribución y/o 
    modificación de este documento bajo los términos de la licencia de
    documentación libre GNU, versión 1.1 o cualquier versión posterior publicada
    por la {\em Free Software Foundation}, sin secciones invariantes. Puede
    consultar esta licencia en http://www.gnu.org. \\[0.2cm]
    Este documento fue compuesto con \LaTeX{}. 
  \end{small}
  \newpage
%   \thispagestyle{empty}
%   \vspace*{0cm}
%   \newpage
}
% -----------------------------------------------------------------------------




% Encabezado y pie de página
\encabezado

\begin{document}

% Silabación extra
\hyphenation{
a-sig-na-tu-ras
au-to-ma-ti-za-rá
ca-li-dad
ca-tá-lo-go
ca-rre-ra
cons-truc-ción
co-rrec-ta-men-te
co-rres-pon-de
de-pen-dien-tes
de-sa-rro-llo
des-tru-yen
diag-nos-tico
dis-po-ni-bi-li-dad
fi-na-li-za-ción
ge-ne-ra-ción
in-fe-rior
man-te-ni-mien-to
me-ca-nis-mo
me-dian-te
o-cu-rren-cia
per-so-nal
pro-ce-di-mien-to
pro-ce-di-mien-tos
pro-por-cio-na-rá
pu-bli-ca-da
re-cha-zar
re-fe-ren-cia
re-lle-nar
re-qui-si-tos
res-pec-ti-va-men-te
res-pecto
res-pon-sa-bi-li-dad
u-su-a-rios
va-ria-rá
vi-lla-rre-al
}


% Portada
\portada{Ingeniería del Software}
{Práctica 2}{Diseño orientado a objetos}
{Sergio de la Rubia García-Carpintero\\Miguel Millán Sánchez-Grande\\
  Luis Muñoz Villarreal\\Alicia Serrano Sánchez\\
  Juan Miguel Torres Triviño}{28 de Mayo de 2010}

% Licencia
\licencia{Sergio de la Rubia García-Carpintero, Miguel Millán Sánchez-Grande,
  Luis Muñoz Villarreal, Alicia Serrano Sánchez, Juan Miguel Torres Triviño}

% Índices
\tableofcontents
%\listoftables
%\listoffigures

\chapter*{Introducción}

\chapter{Especificación de requisitos}

\section{Requisitos iniciales del sistema}
\section{Análisis de requisitos del sistema}

\section{Análisis de requisitos}
\subsection{Usuarios}

\subsection{Paquete de proyectos}

\subsection{Secretario}

\subsection{Coordinador}

\subsection{Área}

\subsection{Adjunto}

\subsection{Subárea}

\subsection{Expertos}

\chapter{Análisis}

\section{Diagrama de casos de uso}

\section{Diagramas de secuencia}

\subsection{Entrada al sistema correcta}

\subsection{Entrada al sistema errónea}

\subsection{Secretario introduce paquete de proyectos}

\subsection{Coordinador selecciona experto para ser adjunto}

\subsection{Coordinador asigna un proyecto a una subárea}

\subsection{Coordinador valida informe}

\subsection{Adjunto busca expertos para evaluar proyecto}

\subsection{Adjunto realiza informe}

\subsection{Experto realiza informe}

\section{Diagrama Entidad/Interrelación}
El modelo entidad/interrelación es el modelo conceptual más utilizado para el 
diseño conceptual de bases de datos.

Es un modelo de red que describe la distribución de los datos almacenados en un 
sistema de forma abstracta. La abstracción busca las propiedades comunes de un 
conjunto de objetos reduciendo la complejidad y ayudando a entender el mundo 
real. Se pretende \emph{visualizar} los objetos que pertenecen a la
\emph{base de datos} como entidades que tienen unos atributos y se vinculan
mediante relaciones.

\imagen{ER.pdf}{14.5}{Diagrama Entidad/Interrelación}{}{}

\chapter{Diseño}

\section{Diagrama de clases}

\section{Diagrama relacional}
Es el modelo más utilizado en la actualidad para modelar problemas reales y 
administrar datos dinámicamente.

Se trata de un modelo lógico, que establece una estructura sobre los datos,
aunque posteriormente éstos puedan ser almacenados de múltiples formas para
aprovechar características físicas concretas de la máquina sobre la que se
implante la base de datos realmente.

\imagen{MRE.pdf}{15}{Diagrama Relacional}{}{}

\chapter{Implementación}

\section{Manual de usuario}

\section{Implementación}

\appendix
\chapter{Carga de trabajo}
\begin{center}
  \begin{tabular}{p{10cm}|c}
    \textbf{Apellidos y Nombre} & \textbf{Porcentaje} \\ \hline \hline
    de la Rubia García-Carpintero, Sergio & 20\% \\
    Millán Sánchez-Grande, Miguel         & 20\% \\ 
    Muñoz Villarreal, Luis                & 20\% \\ 
    Serrano Sánchez, Alicia               & 20\% \\ 
    Torres Triviño, Juan Miguel           & 20\% \\
  \end{tabular}
\end{center}

\end{document}
